% Part: normal-modal-logic
% Chapter: filtrations
% Section: S5-decidable

\documentclass[../../../include/open-logic-section]{subfiles}

\begin{document}

\olfileid[fr]{nml}{fil}{dec}

\olsection{\Log{S5} est décidable}

Pour une classe de modèles dont les modèles finis sont effectivement énumérables,
la propriété du modèle fini, avec une axiomatisation effective, la correction
et la complétude, montre que le système est \emph{décidable} :
il existe une procédure calculable qui détermine si une !!{formula}
est !!{derivable} dans le système.

\begin{thm}
  \Log{S5} est décidable.
\end{thm}

\begin{proof}
  Soit $!A$ donnée, et supposons que les !!{propositional variable}s
  qui figurent dans $!A$ soient parmi $p_1$, \dots, $p_k$. Pour
  chaque $n$, il n'existe qu'un nombre fini de modèles à $n$ mondes
  attribuant des valeurs à $p_1$, \dots, $p_k$. Nous pouvons donc
  mener \emph{en parallèle} deux énumérations : celle de tous les
  théorèmes de \Log{S5}, en engendrant systématiquement les
  démonstrations, et celle des modèles universels finis à $1$, $2$,
  \dots{} mondes, en vérifiant si $!A$ est fausse dans l'un de leurs
  mondes. L'un de ces deux procédés donnera nécessairement une réponse :
  par le \olref[com][fra]{thm:generaldet} et le
  \olref[fmp]{cor:S5fmp}, soit $!A$ est !!{derivable}, soit elle est
  fausse dans un modèle universel fini.
\end{proof}

Cette démonstration fonctionne pour \Log{S5}, car les filtrations de
modèles universels sont toujours universelles. Il en va de même de la
réflexivité et de la sérialité, mais d'autres propriétés demandent
davantage de travail.
Le contre-modèle doit être cherché parmi les modèles de
$\Log{S5}$. Par exemple, la formule $\Box p\lif p$ est
un théorème de $\Log{S5}$, mais elle est fausse dans un
modèle à un monde sans flèche, lorsque $p$ y est fausse.
Les modèles universels finis énumérés dans la preuve
valident tous les théorèmes de~$\Log{S5}$.
\end{document}
